
Behälter

Außer primitiven und zusammengesetzten Daten gibt es
noch Behälter (container).

Unerheblich ist zunächst, um welche Datenstruktur es sich
konkret handelt (Vektor, Menge, Liste, Tabelle, Baum).

Ein einfaches Feld (array) könnte wie folgt notiert werden:

Address[] a = new Address[5];

Das Besondere hierbei ist, dass zwei Typen angegeben werden:
- die eckigen Klammern machen den Wert zum Feld (array);
- der Typ "Address" definiert die darin enthaltenen Elemente.

Doch: Werden hier nicht zwei Ebenen miteinander vermischt?

Betrachtet sei dazu ein Objekt "Person", das mehrere Adressen
enthalten kann:

class Person {
    Address[] a;
}

Bei Verwendung "generischer" Behälter (container) könnte das Beispiel
wie folgt aussehen:

class Person {
    ArrayList<Address> a;
}

Hier hat sich nur die Art des Behälters von "Array" auf "ArrayList"
geändert, doch die enthaltenen Daten sind nach wie vor vom Typ "Address".

Die Klasse "Person" kennt den Typ "ArrayList" ihres Attributes "a",
was einer Meta-Information (Daten über andere Daten) entspricht.
Doch muss sie auch den Typ der in der "ArrayList" enthaltenen
Elemente (also zwei Ebenen tiefer) kennen?

Ist dies nicht eher Aufgabe der "ArrayList"?

... zur besseren Erklärung evtl. doch GUI-Panel mit Array von Knöpfen verwenden?
... dann könnten Position, Farbe etc. als Zusatz-Meta-Information erwähnt werden

Das Problem heutzutage verwendeter Behälter-Datentypen scheint zu sein,
dass sie durch ihre Universalität nicht dazu in der Lage sind,
neben dem Typ noch weitere Meta-Informationen über ihre enthaltenen
Elemente zu speichern.

Betrachtet man das Beispiel in UML-Notation (Abb. xx),
so fällt auf, dass der Behälter-Datentyp einfach ignoriert wird
und beim Attribut nicht auftaucht. Lediglich die Multiplizität
(Kardinalität) * deutet auf die Vielfachheit hin.

ABBILDUNG:

UML-CD mit Klasse "Person" und Attribut a: Address ohne Methoden,
und zweite Klasse "Address" ohne alles; Sternchen an Assoziation von Person zu Address

Zwei mögliche Lösungen sind offensichtlich:

1 Erweiterung der Behälter um Zusatz-Meta-Informationen über ihre enthaltenen Elemente.
Zwar wurde dies in vielen Klassenbibliotheken durch die so genannten "Generics"
versucht, doch können diese nur den Typ der enthaltenen Elemente, nicht aber
weitere Informationen über sie speichern.

2 Erweiterung sämtlicher anderer Typen in der Weise, dass sie alle zum Behälter werden.
Dazu müsste ein Datentyp geschaffen werden, der alle nötigen Behälter-Zugriffsfunktionen
enthielte, zusätzlich aber auch alle relevanten Meta-Informationen über
enthaltene Elemente, wie sie in diesem Artikel untersucht werden.

Ein solcher Datentyp würde keine fest fixierten und übersetzten (compilierten)
Attribute mehr kennen, sondern nur noch mit dynamisch zur Laufzeit
zugewiesenen Elementen arbeiten.

evtl. ABB.?
